\documentclass[11pt]{article}

% ---------- Packages ----------
\usepackage[margin=1in]{geometry}
\usepackage{amsmath, amssymb, amsthm}
\usepackage{mathtools}
\usepackage{enumitem}
\usepackage{graphicx}
\usepackage{booktabs}
\usepackage{array}
\usepackage{hyperref}
\usepackage{xcolor}
\usepackage{fancyhdr}
\usepackage{lastpage}
\usepackage{listings}
\usepackage{courier}

% ---------- Hyperref ----------
\hypersetup{
    colorlinks=true,
    linkcolor=blue,
    urlcolor=blue
}

% ---------- Header/Footer ----------
\pagestyle{fancy}
\fancyhf{}

\lhead{MIT 6.1200J}
\chead{Problem Set 3}
\rhead{Nicolas Iskandar}

\cfoot{\thepage\ /\ \pageref{LastPage}}

% ---------- Custom Commands ----------
\newcommand{\R}{\mathbb{R}}
\newcommand{\N}{\mathbb{N}}
\newcommand{\Z}{\mathbb{Z}}

\newcommand{\answer}{
    \vspace{0.5em}
    \noindent\textbf{Answer.}
}

% ---------- Code Style ----------
\lstset{
    basicstyle=\ttfamily\small,
    breaklines=true,
    frame=single,
    columns=fullflexible,
    keepspaces=true
}

% ---------- Title ----------
\title{
    \vspace{-2em}
    MIT 6.1200J \\
    Problem Set 3
}
\author{Nicolas Iskandar}
\date{}

% ==================================================
\begin{document}

\maketitle

\section*{Problem 1}

\begin{enumerate}[label=\textbf{(\alph*)}]

\item \textbf{States}: The set of all permutations of $\{0, 1, \ldots, n-1\}$.

\textbf{Start state}: The initial permutation $(x_0, x_1, \ldots, x_{n-1})$.

\textbf{Transitions}: From a permutation $(a_0, \ldots, a_{n-1})$, there is a
transition to
\[
(a_0, \ldots, a_{i-1}, a_j, a_{j-1}, \ldots, a_{i+1}, a_i,
a_{j+1}, \ldots, a_{n-1})
\]
for every inverted pair $(i,j)$ with $i<j$ and $a_i>a_j$
(i.e., only good reversals are allowed).

\item \textbf{Final states}: A state has no outgoing transitions when no good
reversal is possible, i.e., when there are no inverted pairs. A permutation has no
inverted pairs if and only if it is sorted in increasing order. Therefore, the unique
final state is the sorted permutation $(0, 1, 2, \ldots, n-1)$.

\textbf{Partial correctness}: If the machine terminates, it must be in a final
state. Since the only final state is $(0, 1, \ldots, n-1)$, the list must be
sorted. $\blacksquare$

\item We define the following derived variable. Interpret the permutation
$(a_0, a_1, \ldots, a_{n-1})$ as the digits of a base-$n$ number:
\[
    V(a_0, \ldots, a_{n-1}) := \sum_{k=0}^{n-1} a_k \cdot n^{n-1-k}.
\]
We claim $V$ is a strictly decreasing derived variable with nonneg integer values.

\textbf{Property (i) --- nonneg integer values}: Since each $a_k \in \{0, \ldots,
n-1\}$ and $n^{n-1-k} \in \mathbb{N}$, we have $V \in \mathbb{N}$. \checkmark

\textbf{Property (ii) --- strictly decreasing}: Suppose we perform a good reversal
on the inverted pair $(i, j)$, where $i < j$ and $a_i > a_j$. The new permutation
$(a_0', \ldots, a_{n-1}')$ satisfies:
\begin{itemize}
    \item $a_k' = a_k$ for $k < i$ and $k > j$ (unchanged outside the reversed
    segment),
    \item the segment $(a_i', \ldots, a_j')$ is the reverse of $(a_i, \ldots, a_j)$,
    so in particular $a_i' = a_j$ and $a_j' = a_i$.
\end{itemize}
% Since only positions $i$ through $j$ change, and each weight $n^{n-1-k}$ is
% divisible by $n^{n-1-j}$ for all $k \leq j$, the difference $V' - V$ is an
% integer multiple of $n^{n-1-j}$. $.

Since $a_k' = a_k$ for $k < i$ and $k > j$, only positions $i$ through $j$
contribute to $V' - V$:
\[
    V' - V = \sum_{k=i}^{j}(a_k' - a_k)\cdot n^{n-1-k}.
\]
Each coefficient $a_k' - a_k$ is an integer, and each weight $n^{n-1-k}$ for
$k \leq j$ satisfies $n - 1 - k \geq n - 1 - j$, so $n^{n-1-j} \mid n^{n-1-k}$.
Therefore every term is divisible by $n^{n-1-j}$, and so $V' - V$ is an integer
multiple of $n^{n-1-j}$. It therefore suffices to show $V' - V < 0$: since $V' - V$
is an integer multiple of $n^{n-1-j}$, a negative value must be at most $-n^{n-1-j}$,
and $n^{n-1-j} \geq 1$ gives us $V' - V \leq -n^{n-1-j} \leq -1$.

To show $V' - V < 0$, factor out $n^{n-1-j}$ from the difference:
\[
    \frac{V' - V}{n^{n-1-j}} = \sum_{k=i}^{j}(a_k' - a_k)\cdot n^{j-k}.
\]
This is the difference between two base-$n$ numbers formed from the same
multiset of digits. In the new number, the most significant position $k = i$
has digit $a_i' = a_j$, while in the old it had digit $a_i > a_j$. Therefore
the most significant digit has strictly decreased, which means the new number
is strictly smaller than the old. Hence $V' - V < 0$, and therefore:
\[
    V' - V \leq -n^{n-1-j} \leq -1 < 0.
\]
So $V$ strictly decreases with every good reversal. \checkmark

\item \textbf{Total correctness}: We combine parts (b) and (c).

\textbf{Termination}: By part (c), $V$ is a derived variable that takes only
nonneg integer values and strictly decreases with every transition. Since a
nonneg integer cannot decrease forever, every execution must eventually
terminate.

\textbf{Partial correctness}: By part (b), if the machine terminates, it must
be in a final state, and the unique final state is the sorted permutation
$(0, 1, \ldots, n-1)$.

Combining these two facts, every execution eventually terminates, and upon
termination the list is fully sorted. $\blacksquare$

\end{enumerate}

\section*{Problem 2}

\begin{enumerate}[label=\textbf{(\alph*)}]

\item We use the perturbation method. Let
\[
    S = \sum_{n=1}^{\infty} \frac{F_n}{c^n}.
\]
Divide both sides by $c$:
\[
    \frac{S}{c} = \sum_{n=1}^{\infty} \frac{F_n}{c^{n+1}} = \sum_{n=2}^{\infty}
    \frac{F_{n-1}}{c^n}.
\]
Similarly,
\[
    \frac{S}{c^2} = \sum_{n=1}^{\infty} \frac{F_n}{c^{n+2}} = \sum_{n=3}^{\infty}
    \frac{F_{n-2}}{c^n}.
\]
Now compute $S - \frac{S}{c} - \frac{S}{c^2}$:
\[
    S - \frac{S}{c} - \frac{S}{c^2} = \sum_{n=1}^{\infty} \frac{F_n}{c^n} -
    \sum_{n=2}^{\infty} \frac{F_{n-1}}{c^n} - \sum_{n=3}^{\infty}
    \frac{F_{n-2}}{c^n}.
\]
Separating the first two terms:
\[
    = \frac{F_1}{c} + \frac{F_2 - F_1}{c^2} + \sum_{n=3}^{\infty}
    \frac{F_n - F_{n-1} - F_{n-2}}{c^n}.
\]
By the Fibonacci recurrence, $F_n - F_{n-1} - F_{n-2} = 0$ for all $n \geq 3$.
Also $F_1 = 1$ and $F_2 = 1$, so:
\[
    S\left(1 - \frac{1}{c} - \frac{1}{c^2}\right) = \frac{1}{c} + \frac{0}{c^2} =
    \frac{1}{c}.
\]
Therefore:
\[
    S = \frac{1/c}{1 - 1/c - 1/c^2} = \frac{1/c \cdot c^2}{c^2 - c - 1} =
    \frac{c}{c^2 - c - 1}.
\]
Substituting $c = 10^5$:
\[
    \boxed{S = \frac{10^5}{10^{10} - 10^5 - 1} = \frac{100000}{9999899999}.}
\]

\item Using WolframAlpha to compute $S = \frac{100000}{9999899999}$ to 50 decimal
places:
\[
    S \approx 0.{\bf 00001}\ 00001\ 00002\ 00003\ 00005\ 00008\ 00013\ 00021\ 00034\
    00055\ldots
\]
The pattern is clear: the digits after the decimal point, read in blocks of 5, give
exactly the Fibonacci numbers $F_1, F_2, F_3, \ldots = 1, 1, 2, 3, 5, 8, 13, 21,
34, 55, \ldots$

This makes sense: by definition $S = \sum_{n=1}^{\infty} F_n / c^n$, so each
Fibonacci number occupies its own block of 5 decimal digits (as long as $F_n < c =
10^5$, i.e., before carrying occurs).

\end{enumerate}

\section*{Problem 3}

Let $f(x) = \dfrac{x}{(5+x)^3}$.

\begin{enumerate}[label=\textbf{(\alph*)}]

\item The integral method requires $f$ to be eventually non-negative, continuous,
and decreasing. While $f(x) > 0$ and is eventually decreasing, the issue is at
$x = 0$: $f(0) = 0$. More critically, $f$ is not monotonically decreasing starting
from $x = 0$. We verify: $f'(x) = \frac{(5+x)^3 - 3x(5+x)^2}{(5+x)^6} =
\frac{5-2x}{(5+x)^4}$, which is positive for $x < 5/2$ and negative for $x > 5/2$.
So $f$ is \textit{increasing} on $[0, 5/2]$, meaning the integral method cannot be
applied directly to $S_0$.

\item We need $f$ to be decreasing on $[i, \infty)$, which requires $i > 5/2$.
Since $i$ must be an integer, the smallest such index is $\boxed{i = 3}$.

\item We apply the integral method to $S_3 = \sum_{x=3}^{\infty} f(x)$. Since $f$
is decreasing on $[3, \infty)$:
\[
    \int_3^{\infty} f(x)\,dx \leq S_3 \leq f(3) + \int_3^{\infty} f(x)\,dx.
\]
We compute the integral. Using the substitution $u = 5 + x$:
\[
    \int_3^{\infty} \frac{x}{(5+x)^3}\,dx = \int_8^{\infty} \frac{u-5}{u^3}\,du =
    \int_8^{\infty} \left(\frac{1}{u^2} - \frac{5}{u^3}\right)du =
    \left[-\frac{1}{u} + \frac{5}{2u^2}\right]_8^{\infty} = \frac{1}{8} -
    \frac{5}{128} = \frac{11}{128}.
\]
Also, $f(3) = \frac{3}{8^3} = \frac{3}{512}$. So:
\[
    \frac{11}{128} \leq S_3 \leq \frac{11}{128} + \frac{3}{512} =
    \frac{44}{512} + \frac{3}{512} = \frac{47}{512}.
\]
Now $S_0 = f(0) + f(1) + f(2) + S_3$. We compute:
\[
    f(0) = 0, \quad f(1) = \frac{1}{216}, \quad f(2) = \frac{2}{343}.
\]
Let $C = f(0) + f(1) + f(2) = \frac{1}{216} + \frac{2}{343}$. Then:
\[
    \frac{11}{128} + C \leq S_0 \leq \frac{47}{512} + C.
\]
The accuracy achieved is $\frac{47}{512} - \frac{11}{128} = \frac{47}{512} -
\frac{44}{512} = \frac{3}{512} \approx 0.00586$.

\item We want accuracy $\frac{1}{200}$, i.e., $f(j) \leq \frac{1}{200}$. Since the
width of the interval from the integral method is $f(j)$ (the upper and lower bounds
differ by exactly $f(j)$), we need the smallest integer $j \geq 3$ such that
$f(j) \leq \frac{1}{200}$.

We check:
\[
    f(j) = \frac{j}{(5+j)^3} \leq \frac{1}{200}.
\]
Testing $j = 7$: $f(7) = \frac{7}{12^3} = \frac{7}{1728} < \frac{1}{200}$ since
$7 \cdot 200 = 1400 < 1728$. \checkmark

Testing $j = 6$: $f(6) = \frac{6}{11^3} = \frac{6}{1331}$. $\frac{6}{1331} \leq
\frac{1}{200}$? We need $6 \cdot 200 \leq 1331$, i.e., $1200 \leq 1331$. \checkmark

Testing $j = 5$: $f(5) = \frac{5}{10^3} = \frac{5}{1000} = \frac{1}{200}$. We need
$\frac{1}{200} \leq \frac{1}{200}$. \checkmark

Testing $j = 4$: $f(4) = \frac{4}{9^3} = \frac{4}{729}$. Is $\frac{4}{729} \leq
\frac{1}{200}$? We need $4 \cdot 200 \leq 729$, i.e., $800 \leq 729$. \texttimes

So the smallest such $j$ is $\boxed{j = 5}$.

Now we apply the integral method to $S_5$:
\[
    \int_5^{\infty} \frac{x}{(5+x)^3}\,dx = \int_{10}^{\infty}
    \frac{u-5}{u^3}\,du = \left[-\frac{1}{u} + \frac{5}{2u^2}\right]_{10}^{\infty}
    = \frac{1}{10} - \frac{5}{200} = \frac{1}{10} - \frac{1}{40} = \frac{3}{40}.
\]
So:
\[
    \frac{3}{40} \leq S_5 \leq \frac{3}{40} + f(5) = \frac{3}{40} + \frac{1}{200} =
    \frac{15}{200} + \frac{1}{200} = \frac{16}{200} = \frac{2}{25}.
\]
Now $S_0 = f(0) + f(1) + f(2) + f(3) + f(4) + S_5$. Computing:
\[
    f(0)=0,\quad f(1)=\frac{1}{216},\quad f(2)=\frac{2}{343},\quad
    f(3)=\frac{3}{512},\quad f(4)=\frac{4}{729}.
\]
Let $C' = \frac{1}{216} + \frac{2}{343} + \frac{3}{512} + \frac{4}{729}$. Then:
\[
    u = \frac{3}{40} + C' \leq S_0 \leq \frac{2}{25} + C' = v,
\]
with $|v - u| = \frac{1}{200}$. \checkmark

\end{enumerate}

\newpage

\section*{Problem 4}

\begin{enumerate}[label=\textbf{(\alph*)}]

\item We write:
\[
    \sum_{i=0}^{n} \frac{5^i - 7^i}{13^i} = \sum_{i=0}^{n} \left(\frac{5}{13}
    \right)^i - \sum_{i=0}^{n}\left(\frac{7}{13}\right)^i.
\]
Using the geometric series formula $\sum_{i=0}^{n} r^i = \frac{1 - r^{n+1}}{1-r}$:
\[
    = \frac{1 - (5/13)^{n+1}}{1 - 5/13} - \frac{1-(7/13)^{n+1}}{1-7/13} =
    \frac{13}{8}\left(1 - \frac{5^{n+1}}{13^{n+1}}\right) - \frac{13}{6}\left(1 -
    \frac{7^{n+1}}{13^{n+1}}\right).
\]
Therefore:
\[
    \boxed{\sum_{i=0}^{n} \frac{5^i - 7^i}{13^i} = \frac{13}{8} - \frac{13}{6} +
    \frac{13}{6} \cdot \frac{7^{n+1}}{13^{n+1}} - \frac{13}{8} \cdot
    \frac{5^{n+1}}{13^{n+1}} = -\frac{13}{24} + \frac{7^{n+1}}{6 \cdot 13^n} -
    \frac{5^{n+1}}{8 \cdot 13^n}.}
\]

\item We evaluate the double sum by switching the order of summation:
\[
    \sum_{i=1}^{n}\sum_{j=1}^{k}(3j - i).
\]
We split and evaluate each part separately:
\[
    \sum_{i=1}^{n}\sum_{j=1}^{k}(3j-i) = \sum_{i=1}^{n}\left(3\sum_{j=1}^{k}j -
    \sum_{j=1}^{k}i\right) = \sum_{i=1}^{n}\left(\frac{3k(k+1)}{2} - ki\right).
\]
Now summing over $i$:
\[
    = n \cdot \frac{3k(k+1)}{2} - k\sum_{i=1}^{n}i = \frac{3nk(k+1)}{2} -
    k\cdot\frac{n(n+1)}{2}.
\]
Therefore:
\[
    \boxed{\sum_{i=1}^{n}\sum_{j=1}^{k}(3j-i) = \frac{nk}{2}\Big(3(k+1) -
    (n+1)\Big) = \frac{nk(3k - n + 2)}{2}.}
\]

\end{enumerate}

\end{document}